\medskip

On considère une fonction $f$ définie sur l'intervalle $]0~;~ +\infty[$. On admet~ qu'elle est deux fois dérivable sur l'intervalle $]0~;~ +\infty[$. On note $f'$ sa fonction dérivée et $f''$ sa
fonction dérivée seconde.

\smallskip

Dans un repère orthogonal, on a tracé ci-dessous:

\begin{itemize}[label=$\bullet~$]
\item la courbe représentative de $f$, notée $\mathcal{C}_f$ sur l'intervalle ]0~;~3] ; 
\item la droite $T_{\text{A}}$,tangente à $\mathcal{C}_f$ au point A(1~;~2) ;
\item la droite $T_{\text{B}}$ tangente à $\mathcal{C}_f$ au point B(e~;~e).
\end{itemize}

\smallskip

On précise par ailleurs que la tangente $T_{\text{A}}$ passe par le point C(3~;~0).

\begin{center}
\psset{xunit=3.75cm,yunit=3.75cm,arrowsize=2pt 3}
\begin{pspicture*}(-0.2,-1.05)(3.3,3.3)
\uput[ur](1,2){A}\uput[ul](2.71828,2.71828){B}\uput[ur](3,0){C}\uput[ur](-0.2,3.15){$T_{\text{A}}$}
\uput[d](2.71828,0){e}\uput[l](0,2.71828){e}\uput[d](0.2,0){0,2}
\uput[dr](3,3.2){$T_{\text{B}}$}\uput[u](0.2,2.45){\red $\mathcal{C}_f$}
\multido{\n=0.0+0.2}{17}{\psline[linewidth=0.25pt](\n,-1)(\n,3.3)}
\multido{\n=-1+1}{5}{\psline[linewidth=0.25pt](0,\n)(3.3,\n)}
\psaxes[linewidth=1.25pt]{->}(0,0)(0,0)(3.3,3.3)
\psaxes[linewidth=1.25pt](0,0)(0,0)(1,1)

\psplotTangent[linewidth=1pt]{1}{3}{x ln dup mul 2 mul x ln 3 mul sub 2 add x mul}
\psplotTangent[linewidth=1pt]{2.71828}{4}{x ln dup mul 2 mul x ln 3 mul sub 2 add x mul}
\psdots(1,2)(2.71828,2.71828)
\psline[linestyle=dashed,linewidth=1.5pt](0,2.7128)(2.71828,2.71828)(2.71828,0)
\pscustom[fillstyle=vlines,hatchsep=8pt,hatchcolor=blue,linecolor=blue]{(\psline(2.71828,2.71828)(1,-0.718)(1,2)\psplot[plotpoints=2000,linewidth=1.25pt,linecolor=blue]{1}{2.71828}{x ln dup mul 2 mul x ln 3 mul sub 2 add x mul}}
\psplot[plotpoints=2000,linewidth=1.35pt,linecolor=red]{0.0025}{3.3}{x ln dup mul 2 mul x ln 3 mul sub 2 add x mul}
\end{pspicture*}
\end{center}

\begin{center}\textbf{Partie A :} Lectures graphiques\end{center}

On répondra aux questions suivantes en les justifiant à l'aide du graphique.

\medskip
\begin{enumerate}
\item Déterminer le nombre dérivé $f'(1)$.
\item Combien de solutions l'équation $f'(x) = 0$ admet-elle dans l'intervalle ]0~;~3] ?
\item Quel est le signe de $f''(0,2)$ ?
\end{enumerate}


\begin{center}\textbf{Partie B :} étude de la fonction $f$\end{center}

On admet dans cette partie que la fonction $f$ est définie sur l'intervalle $]0~;~ +\infty[$ par

\[f(x) = x[2(\ln x)^2 - 3\ln x + 2 ]\]

où $\ln$ désigne la fonction logarithme népérien.

\medskip

\begin{enumerate}
\item Résoudre dans $\R$ l'équation $2X^2- 3X + 2 = 0$.

En déduire que $\mathcal{C}_f$ ne coupe pas l'axe des abscisses.
\item Déterminer, en justifiant, la limite de $f$ en $+\infty$.

On admettra que la limite de $f$ en 0 est égale à 0.
\item On admet que pour tout $x$ appartenant à $]0~;~ +\infty[$, \: $f'(x) = 2(\ln x)^2 + \ln x - 1$.
	\begin{enumerate}
		\item Montrer que pour tout $x$ appartenant à $]0~;~ +\infty[$,\: $f''(x)= \dfrac 1x(4\ln x + 1)$.
		\item Étudier la convexité de la fonction $f$ sur l'intervalle $]0~;~ +\infty[$ et préciser la valeur exacte de l'abscisse du point d'inflexion.
		\item Montrer que la courbe $\mathcal{C}_f$ est au-dessus de la tangente $T_{\text{B}}$ sur l'intervalle $]1~;~ +\infty[$.
	\end{enumerate}
\end{enumerate}

\begin{center}\textbf{Partie C :} Calcul d'aire\end{center}

\begin{enumerate}
\item Justifier que la tangente $T_{\text{B}}$ a pour équation réduite $y = 2x - \e$.


\item À l'aide d'une intégration par parties, montrer que

\[\displaystyle\int_1^{\e} x\ln x\, \text{d}x = \dfrac{\e^2 + 1}{4}.\]
\item On note $\mathcal{A}$ l'aire du domaine hachuré sur la figure, délimité par la courbe $\mathcal{C}_f$, la tangente $T_{\text{B}}$, et les droites d'équation $x = 1$ et $x = \e$.

On admet que $\displaystyle\int_1^{\e} x(\ln x)^2\,\text{d}x = \dfrac{\e^2 - 1}{4}$.

En déduire la valeur exacte de $\mathcal{A}$ en unité d'aire.
\end{enumerate}

\bigskip
